\chapter{Business Intelligence et tableaux de bord operationnels}

\section{Introduction}
% TODO: Introduire le role de Power BI dans le pilotage operationnel.

\section{Objectifs des dashboards}
% TODO: Expliquer les objectifs management, boss et field team.

\section{Definition des KPI operationnels}
% TODO: Decrire les KPI reels utilises sans inventer de resultats.

\section{Dashboard executif}
\begin{figure}[htbp]
    \centering
    \fbox{\parbox{0.82\textwidth}{\centering Espace reserve a la capture du dashboard executif Power BI.}}
    \caption{Dashboard executif de suivi operationnel}
    \label{fig:dashboard_executif}
\end{figure}

\section{Dashboard Merchandiser Performance Intelligence}
\begin{figure}[htbp]
    \centering
    \fbox{\parbox{0.82\textwidth}{\centering Espace reserve a la capture du dashboard Merchandiser Performance Intelligence.}}
    \caption{Dashboard d'analyse de performance des merchandisers}
    \label{fig:dashboard_merch}
\end{figure}

\section{Analyse GT / MT}
% TODO: Presenter l'analyse GT/MT et son interet business.

\section{Apport pour le management et le field team}
% TODO: Expliquer l'apport des dashboards pour les decisions operationnelles.

\section{Conclusion}
% TODO: Faire la transition vers l'application mobile client.

